% يجب الترجمة بمحرك LuaLaTeX  (Menu -> Compiler -> LuaLaTeX)
\documentclass[12pt, a4paper]{article}
\usepackage[a4paper, top=1.2cm, bottom=1.4cm, left=1.4cm, right=1.4cm]{geometry}
\usepackage{amsmath, amssymb}
\usepackage{fontspec}
\usepackage[bidi=basic]{babel}
\babelprovide[main, import, onchar=ids fonts]{arabic}
\babelprovide[import, onchar=ids fonts]{english}
\babelfont{rm}{Noto Serif}
\babelfont[arabic]{rm}{Amiri}
\usepackage{enumitem}
\usepackage{tabularx}
\usepackage{array}
\usepackage{graphicx}
\usepackage{tikz}
\usepackage{eso-pic}

\setlength{\parindent}{0pt}
\newcommand{\LTRfig}[1]{\mbox{\foreignlanguage{english}{#1}}}

% عنوان سؤال داخل شريط رمادي مع العلامة على اليسار
\newcommand{\question}[2]{%
  \par\vspace{0.25cm}\noindent
  \colorbox{gray!15}{\parbox{\dimexpr\linewidth-2\fboxsep\relax}{\textbf{\large #1}\hfill\mbox{\textbf{(#2)}}}}%
  \par\vspace{0.15cm}}
% علامة الفرع في نهاية السطر
\newcommand{\pts}[1]{\hfill\mbox{\textbf{(#1)}}}
% فرع داخل عمود (للفروع المتجاورة)
\newcommand{\sub}[3]{% #1 الحرف، #2 النص، #3 مساحة الحل
  \textbf{#1)} #2\par\vspace{#3}}

\begin{document}

% إطار الصفحات
\AddToShipoutPictureBG{%
  \LTRfig{%
  \begin{tikzpicture}[remember picture, overlay]
    \draw[line width=1.2pt] ([shift={(0.7cm,-0.7cm)}]current page.north west) rectangle ([shift={(-0.7cm,0.7cm)}]current page.south east);
  \end{tikzpicture}}%
}

% ============================ الترويسة ============================
\noindent
\begin{minipage}[t]{0.37\textwidth}
\vspace{0pt}
\textbf{المملكة الأردنية الهاشمية} \\
\textbf{وزارة التربية والتعليم} \\
\textbf{مدرسة عنبة الثانوية الشاملة للبنات} \\
\textbf{العام الدراسي: 2025 / 2026}
\end{minipage}%
\begin{minipage}[t]{0.26\textwidth}
\vspace{0pt}
\centering
\IfFileExists{logo.png}{\includegraphics[width=2.2cm]{logo.png}}{}
\end{minipage}%
\begin{minipage}[t]{0.37\textwidth}
\vspace{0pt}
\textbf{المبحث: الرياضيات} \\
\textbf{الصف: السابع} \\
\textbf{الزمن: حصة واحدة} \\
\textbf{الدرجة الكلية: 33 علامة}
\end{minipage}

\vspace{0.2cm}
\hrule height 1.2pt
\vspace{0.06cm}
\hrule height 0.4pt
\vspace{0.2cm}

\begin{center}
\textbf{\LARGE اختبار منتصف الفصل الأول}
\end{center}
\vspace{-0.1cm}

% بيانات الطالبة
\noindent\fbox{\parbox{\dimexpr\linewidth-2\fboxsep-2\fboxrule\relax}{%
\vspace{0.15cm}
\textbf{الاسم:} \rule{6.5cm}{0.4pt} \hfill
\textbf{الشعبة:} \rule{2.2cm}{0.4pt} \hfill
\textbf{التاريخ:} \rule{3cm}{0.4pt}
\vspace{0.1cm}}}

\vspace{0.2cm}
% جدول توزيع العلامات
\begin{center}
\renewcommand{\arraystretch}{1.25}
\begin{tabular}{|>{\bfseries}c|*{6}{>{\centering\arraybackslash}p{1.3cm}|}>{\centering\arraybackslash}p{1.6cm}|}
\hline
السؤال & الأول & الثاني & الثالث & الرابع & الخامس & السادس & المجموع \\ \hline
العلامة & 6 & 6 & 6 & 6 & 6 & 3 & \textbf{33} \\ \hline
علامة الطالبة & & & & & & & \\ \hline
\end{tabular}
\end{center}

\vspace{-0.1cm}
\noindent\textbf{تعليمات:} أجب عن جميع الأسئلة، واكتب خطوات الحل في المكان المخصص. الآلة الحاسبة غير مسموحة.

% ============================ السؤال الأول ============================
\question{السؤال الأول: تمثيل الأعداد النسبية بالصورة العشرية والكسرية}{6 علامات}

\noindent
\begin{minipage}[t]{0.32\linewidth}
\sub{أ}{اكتب الكسر $\dfrac{3}{8}$ على الصورة العشرية. \pts{2}}{2.3cm}
\end{minipage}\hfill\vrule\hfill
\begin{minipage}[t]{0.32\linewidth}
\sub{ب}{اكتب العدد العشري $0.45$ على صورة كسر في أبسط صورة. \pts{2}}{2.3cm}
\end{minipage}\hfill\vrule\hfill
\begin{minipage}[t]{0.32\linewidth}
\sub{ج}{اكتب الكسر $\dfrac{4}{9}$ على صورة عدد عشري دوري. \pts{2}}{2.3cm}
\end{minipage}

% ============================ السؤال الثاني ============================
\question{السؤال الثاني: ترتيب الأعداد النسبية ومقارنتها وجمعها}{6 علامات}

\textbf{أ)} رتّب الأعداد النسبية الآتية تصاعديًا: \quad
$-\dfrac{5}{6} \ ,\ -0.8 \ ,\ 65\% \ ,\ \left|-\dfrac{3}{4}\right|$ \pts{2}
\vspace{1.6cm}

\noindent
\begin{minipage}[t]{0.48\linewidth}
\sub{ب}{ضع إشارة $<$ أو $>$ أو $=$ في المربع: \quad $-\dfrac{5}{8} \ \square\ -\dfrac{2}{7}$ \pts{2}}{1.8cm}
\end{minipage}\hfill\vrule\hfill
\begin{minipage}[t]{0.48\linewidth}
\sub{ج}{أوجد ناتج الجمع في أبسط صورة: \quad $-\dfrac{5}{6} + \dfrac{1}{4}$ \pts{2}}{1.8cm}
\end{minipage}

% ============================ السؤال الثالث ============================
\question{السؤال الثالث: ضرب الأعداد النسبية وقسمتها}{6 علامات}

\noindent
\begin{minipage}[t]{0.48\linewidth}
\sub{أ}{أوجد ناتج الضرب في أبسط صورة: \par\vspace{0.1cm}
\centering $-2\dfrac{1}{2} \times \left(-\dfrac{4}{5}\right)$ \pts{3}}{2.9cm}
\end{minipage}\hfill\vrule\hfill
\begin{minipage}[t]{0.48\linewidth}
\sub{ب}{أوجد ناتج القسمة في أبسط صورة: \par\vspace{0.1cm}
\centering $6 \div \left(-2\dfrac{1}{4}\right)$ \pts{3}}{2.9cm}
\end{minipage}

\newpage
% ============================ السؤال الرابع ============================
\question{السؤال الرابع: الأسس وترتيب العمليات الحسابية}{6 علامات}

\noindent
\begin{minipage}[t]{0.32\linewidth}
\sub{أ}{اكتب المقدار الآتي بالصيغة الأسية: \par\vspace{0.1cm}
\centering $\dfrac{1}{2} \times \dfrac{1}{2} \times 3 \times 3 \times 3$ \pts{2}}{2.3cm}
\end{minipage}\hfill\vrule\hfill
\begin{minipage}[t]{0.32\linewidth}
\sub{ب}{أوجد قيمة المقدار: \par\vspace{0.1cm}
\centering $4^{-2}$ \pts{2}}{2.3cm}
\end{minipage}\hfill\vrule\hfill
\begin{minipage}[t]{0.32\linewidth}
\sub{ج}{أوجد قيمة المقدار مع مراعاة أولويات العمليات: \par\vspace{0.1cm}
\centering $3(4-2)^3 + 8 \div 2^2$ \pts{2}}{2.3cm}
\end{minipage}

% ============================ السؤال الخامس ============================
\question{السؤال الخامس: الجبر والتفكير العالي}{6 علامات}

\textbf{أ)} في المقدار الجبري $5x - 3y + 8$ ، حدّد: \pts{2} \par\vspace{0.2cm}
\noindent\begin{tabularx}{\linewidth}{XXX}
الحدود الجبرية: \rule{2.6cm}{0.4pt} & معامل $y$: \rule{2.6cm}{0.4pt} & الحد الثابت: \rule{2.6cm}{0.4pt}
\end{tabularx}
\vspace{0.3cm}

\textbf{ب)} خزان سيارة فيه كمية من الوقود. استهلكت السيارة $6.3\ \mathrm{L}$، ثم أُضيف إليها $15\ \mathrm{L}$، ثم استهلكت $11\dfrac{4}{5}\ \mathrm{L}$، فبقي في الخزان $8.9\ \mathrm{L}$. أوجد الكمية الأصلية من الوقود باستخدام الحل العكسي. \pts{2}
\vspace{2.8cm}

\textbf{ج)} إذا علمت أن الأعداد النسبية $a < b < c < 0$ ، كيف يمكن ترتيب معكوساتها الضربية
$\dfrac{1}{a} \ ,\ \dfrac{1}{b} \ ,\ \dfrac{1}{c}$ ؟ برّر إجابتك بمثال عددي. \pts{2}
\vspace{2.8cm}

% ============================ السؤال السادس ============================
\question{السؤال السادس: صح أم خطأ}{3 علامات}

اكتب \textbf{T} داخل القوس إذا كانت الجملة صحيحة، واكتب \textbf{F} داخل القوس إذا كانت الجملة خاطئة. \hfill (0.5 علامة لكل جملة)

\begin{enumerate}[label=\textbf{\arabic*)}, itemsep=0.12cm, leftmargin=*]
\item \textbf{( \hspace{1cm} )} \ كل عدد عشري منتهٍ يمكن كتابته على صورة كسر مقامه قوة للعدد $10$.
\item \textbf{( \hspace{1cm} )} \ كل عدد عشري دوري ليس عددًا نسبيًا.
\item \textbf{( \hspace{1cm} )} \ إذا كان $a < 0$ و $b < 0$ و $a > b$ فإن $|a| < |b|$.
\item \textbf{( \hspace{1cm} )} \ المعكوس الضربي للعدد $-\dfrac{3}{4}$ هو $\dfrac{4}{3}$.
\item \textbf{( \hspace{1cm} )} \ ناتج طرح عدد سالب من عدد سالب آخر يكون دائمًا عددًا سالبًا.
\item \textbf{( \hspace{1cm} )} \ قيمة العدد $2^{-3}$ تساوي $\dfrac{1}{8}$.
\end{enumerate}

\vfill
\begin{center}
\textbf{\large انتهت الأسئلة، مع تمنياتي لكم بالتوفيق}
\end{center}

\end{document}
